\documentclass[UTF8,11pt]{ctexart}
\usepackage[a4paper,margin=25mm]{geometry}
\usepackage{amsmath,amssymb,amsthm,mathtools}
\usepackage[all]{xy}
\usepackage{xurl}
\usepackage{hyperref}
\usepackage{enumitem}
\hypersetup{hidelinks,breaklinks=true}
\setlength{\emergencystretch}{3em}
\newtheorem*{theorem}{Theorem}
\newtheorem*{lemma}{Lemma}
\newtheorem*{proposition}{Proposition}
\newtheorem*{corollary}{Corollary}
\newtheorem*{definition}{Definition}
\newcommand{\Spec}{\operatorname{Spec}}
\newcommand{\Hom}{\operatorname{Hom}}
\newcommand{\End}{\operatorname{End}}
\newcommand{\Gal}{\operatorname{Gal}}
\newcommand{\Cl}{\operatorname{Cl}}
\newcommand{\vol}{\operatorname{vol}}
\newcommand{\covol}{\operatorname{covol}}
\newcommand{\tr}{\operatorname{tr}}
\newcommand{\im}{\operatorname{im}}
\newcommand{\id}{\operatorname{id}}
\newcommand{\coker}{\operatorname{coker}}
\newcommand{\stacksref}[2]{\href{https://stacks.math.columbia.edu/tag/#1}{#2 [#1]}}
\begin{document}
\section*{006\quad Auslander--Buchsbaum深度公式}
\subsection*{来源与整理范围}
The Stacks Project Authors, \emph{The Stacks Project}, Chapter 10, \emph{Commutative Algebra}。主命题：Proposition 10.111.1，Auslander--Buchsbaum formula，Tag 090V。收录其完整证明，并附原文用于关键深度下界的 Proposition 10.102.9（00N1）及两个支撑引理 10.102.2（00MT）、10.102.8（00N0）。这些支撑材料只服务于本题，不另计题数。
\par\url{https://stacks.math.columbia.edu/tag/090V}
\par\url{https://stacks.math.columbia.edu/tag/00N1}
\par\url{https://stacks.math.columbia.edu/tag/00MT}
\par\url{https://stacks.math.columbia.edu/tag/00N0}
\par 记号出处：Situation 10.102.1（\url{https://stacks.math.columbia.edu/tag/00MS}）和 Definition 10.102.5（\url{https://stacks.math.columbia.edu/tag/00MV}）。使用版本：2026-09-28 实际访问的在线正文；未取得该网页构建的固定 commit。以下正文由网页中的数学 TeX 整理，保留作者论证及原编号引用。

\subsection*{作者原命题与人类证明}
\paragraph{Situation 10.102.1.}
Here $R$ is a ring, and we have a complex
\[
0\to R^{n_e}\xrightarrow{\varphi_e}R^{n_{e-1}}\to\ldots
\xrightarrow{\varphi_{i+1}}R^{n_i}\xrightarrow{\varphi_i}R^{n_{i-1}}
\to\ldots\xrightarrow{\varphi_1}R^{n_0}.
\]
In other words we require $\varphi_i\circ\varphi_{i+1}=0$ for $i=1,\ldots,e-1$.

\begin{definition}[10.102.5, Tag 00MV]
Let $R$ be a ring. Suppose that $\varphi:R^m\to R^n$ is a map of finite free $R$-modules.
\begin{enumerate}
\item The rank of $\varphi$ is the maximal $r$ such that $\wedge^r\varphi:\wedge^rR^m\to\wedge^rR^n$ is nonzero.
\item We denote $I(\varphi)\subset R$ the ideal generated by the $r\times r$ minors of the matrix of $\varphi$, where $r$ is the rank of $\varphi$ as defined above.
\end{enumerate}
\end{definition}

\begin{lemma}[10.102.2, Tag 00MT]
Suppose $R$ is a ring. Let
\[
\ldots\xrightarrow{\varphi_{i+1}}R^{n_i}\xrightarrow{\varphi_i}R^{n_{i-1}}
\xrightarrow{\varphi_{i-1}}\ldots
\]
be a complex of finite free $R$-modules. Suppose that for some $i$ some matrix coefficient of the map $\varphi_i$ is invertible. Then the displayed complex is isomorphic to the direct sum of a complex
\[
\begin{aligned}
\ldots\to R^{n_{i+2}}\xrightarrow{\varphi_{i+2}}R^{n_{i+1}}
&\to R^{n_i-1}\to R^{n_{i-1}-1}\\
&\to R^{n_{i-2}}\xrightarrow{\varphi_{i-2}}R^{n_{i-3}}\to\ldots
\end{aligned}
\]
and the complex $\ldots\to0\to R\to R\to0\to\ldots$ where the map $R\to R$ is the identity map.
\end{lemma}
\begin{proof}
The assumption means, after a change of basis of $R^{n_i}$ and $R^{n_{i-1}}$ that the first basis vector of $R^{n_i}$ is mapped via $\varphi_i$ to the first basis vector of $R^{n_{i-1}}$. Let $e_j$ denote the $j$th basis vector of $R^{n_i}$ and $f_k$ the $k$th basis vector of $R^{n_{i-1}}$. Write $\varphi_i(e_j)=\sum a_{jk}f_k$. So $a_{1k}=0$ unless $k=1$ and $a_{11}=1$. Change basis on $R^{n_i}$ again by setting $e'_j=e_j-a_{j1}e_1$ for $j>1$. After this change of coordinates we have $a_{j1}=0$ for $j>1$. Note the image of $R^{n_{i+1}}\to R^{n_i}$ is contained in the subspace spanned by $e_j$, $j>1$. Note also that $R^{n_{i-1}}\to R^{n_{i-2}}$ has to annihilate $f_1$ since it is in the image. These conditions and the shape of the matrix $(a_{jk})$ for $\varphi_i$ imply the lemma.
\end{proof}

\begin{lemma}[10.102.8 (Acyclicity lemma), Tag 00N0]
Let $R$ be a local Noetherian ring. Let $0\to M_e\to M_{e-1}\to\ldots\to M_0$ be a complex of finite $R$-modules. Assume $\operatorname{depth}(M_i)\geq i$. Let $i$ be the largest index such that the complex is not exact at $M_i$. If $i>0$ then $\operatorname{Ker}(M_i\to M_{i-1})/\operatorname{Im}(M_{i+1}\to M_i)$ has depth $\geq1$.
\end{lemma}
\begin{proof}
Let $H=\operatorname{Ker}(M_i\to M_{i-1})/\operatorname{Im}(M_{i+1}\to M_i)$ be the cohomology group in question. We may break the complex into short exact sequences
\[
\begin{gathered}
0\to M_e\to M_{e-1}\to K_{e-2}\to0,\\
0\to K_j\to M_j\to K_{j-1}\to0,\qquad i+2\leq j\leq e-2,\\
0\to K_{i+1}\to M_{i+1}\to B_i\to0,\\
0\to K_i\to M_i\to M_{i-1},
\end{gathered}
\]
and $0\to B_i\to K_i\to H\to0$. We proceed up through these complexes to prove the statements about depths, repeatedly using Lemma 10.72.6. First of all, since $\operatorname{depth}(M_e)\geq e$, and $\operatorname{depth}(M_{e-1})\geq e-1$ we deduce that $\operatorname{depth}(K_{e-2})\geq e-1$. At this point the sequences $0\to K_j\to M_j\to K_{j-1}\to0$ for $i+2\leq j\leq e-2$ imply similarly that $\operatorname{depth}(K_{j-1})\geq j$ for $i+2\leq j\leq e-2$. The sequence $0\to K_{i+1}\to M_{i+1}\to B_i\to0$ then shows that $\operatorname{depth}(B_i)\geq i+1$. The sequence $0\to K_i\to M_i\to M_{i-1}$ shows that $\operatorname{depth}(K_i)\geq1$ since $M_i$ has depth $\geq i\geq1$ by assumption. The sequence $0\to B_i\to K_i\to H\to0$ then implies the result.
\end{proof}

\begin{proposition}[10.102.9, Tag 00N1]
In Situation 10.102.1, suppose $R$ is a local Noetherian ring. The following are equivalent
\begin{enumerate}
\item $0\to R^{n_e}\to R^{n_{e-1}}\to\ldots\to R^{n_0}$ is exact at $R^{n_e},\ldots,R^{n_1}$, and
\item for all $i$, $1\leq i\leq e$ the following two conditions are satisfied:
\begin{enumerate}[label=(\alph*)]
\item $\operatorname{rank}(\varphi_i)=r_i$ where $r_i=n_i-n_{i+1}+\ldots+(-1)^{e-i-1}n_{e-1}+(-1)^{e-i}n_e$,
\item $I(\varphi_i)=R$, or $I(\varphi_i)$ contains a regular sequence of length $i$.
\end{enumerate}
\end{enumerate}
\end{proposition}
\begin{proof}
If for some $i$ some matrix coefficient of $\varphi_i$ is not in $\mathfrak m$, then we apply Lemma 10.102.2. It is easy to see that the proposition for a complex and for the same complex with a trivial complex added to it are equivalent. Thus we may assume that all matrix entries of each $\varphi_i$ are elements of the maximal ideal. We may also assume that $e\geq1$.

Assume the complex is exact at $R^{n_e},\ldots,R^{n_1}$. Let $\mathfrak q\in\operatorname{Ass}(R)$. Note that the ring $R_{\mathfrak q}$ has depth $0$ and that the complex remains exact after localization at $\mathfrak q$. We apply Lemmas 10.102.3 and 10.102.6 to the localized complex over $R_{\mathfrak q}$. We conclude that $\varphi_{i,\mathfrak q}$ has rank $r_i$ for all $i$. Since $R\to\bigoplus_{\mathfrak q\in\operatorname{Ass}(R)}R_{\mathfrak q}$ is injective (Lemma 10.63.19), we conclude that $\varphi_i$ has rank $r_i$ over $R$ by the definition of rank as given in Definition 10.102.5. Therefore we see that $I(\varphi_i)_{\mathfrak q}=I(\varphi_{i,\mathfrak q})$ as the ranks do not change. Since all of the ideals $I(\varphi_i)_{\mathfrak q}$, $e\geq i\geq1$ are equal to $R_{\mathfrak q}$ (by the lemmas referenced above) we conclude none of the ideals $I(\varphi_i)$ is contained in $\mathfrak q$. This implies that $I(\varphi_e)I(\varphi_{e-1})\ldots I(\varphi_1)$ is not contained in any of the associated primes of $R$. By Lemma 10.15.2 we may choose $x\in I(\varphi_e)I(\varphi_{e-1})\ldots I(\varphi_1)$, $x\notin\mathfrak q$ for all $\mathfrak q\in\operatorname{Ass}(R)$. Observe that $x$ is a nonzerodivisor (Lemma 10.63.9). According to Lemma 10.102.7 the complex $0\to(R/xR)^{n_e}\to\ldots\to(R/xR)^{n_1}$ is exact at $(R/xR)^{n_e},\ldots,(R/xR)^{n_2}$. By induction on $e$ all the ideals $I(\varphi_i)/xR$ have a regular sequence of length $i-1$. This proves that $I(\varphi_i)$ contains a regular sequence of length $i$.

Assume (2)(a) and (2)(b) hold. We will prove that (1) holds by induction on $\dim(R)$. If $\dim(R)=0$, then we must have $I(\varphi_i)=R$ for $1\leq i\leq e$ by (2)(b). Since the coefficients of $\varphi_i$ are contained in the maximal ideal this can happen only if $r_i=0$ for all $i$. By (2)(a) we conclude that $e=0$ and (1) holds. Assume $\dim(R)>0$. We claim that for any prime $\mathfrak p\subset R$ conditions (2)(a) and (2)(b) hold for the complex $0\to R_{\mathfrak p}^{n_e}\to R_{\mathfrak p}^{n_{e-1}}\to\ldots\to R_{\mathfrak p}^{n_0}$ with maps $\varphi_{i,\mathfrak p}$ over $R_{\mathfrak p}$. Namely, since $I(\varphi_i)$ contains a nonzero divisor, the image of $I(\varphi_i)$ in $R_{\mathfrak p}$ is nonzero. This implies that the rank of $\varphi_{i,\mathfrak p}$ is the same as the rank of $\varphi_i$: the rank as defined above of a matrix $\varphi$ over a ring $R$ can only drop when passing to an $R$-algebra $R'$ and this happens if and only if $I(\varphi)$ maps to zero in $R'$. Thus (2)(a) holds. Having said this we know that $I(\varphi_{i,\mathfrak p})=I(\varphi_i)_{\mathfrak p}$ and we see that (2)(b) is preserved under localization as well. By induction on the dimension of $R$ we may assume the complex is exact when localized at any nonmaximal prime $\mathfrak p$ of $R$. Thus $\operatorname{Ker}(\varphi_i)/\operatorname{Im}(\varphi_{i+1})$ has support contained in $\{\mathfrak m\}$ and hence if nonzero has depth $0$. As $I(\varphi_i)\subset\mathfrak m$ for all $i$ because of what was said in the first paragraph of the proof, we see that (2)(b) implies $\operatorname{depth}(R)\geq e$. By Lemma 10.102.8 we see that the complex is exact at $R^{n_e},\ldots,R^{n_1}$ concluding the proof.
\end{proof}

\begin{proposition}[10.111.1 (Auslander--Buchsbaum), Tag 090V]
Let $R$ be a Noetherian local ring. Let $M$ be a nonzero finite $R$-module which has finite projective dimension $\operatorname{pd}_R(M)$. Then we have
\[
\operatorname{depth}(R)=\operatorname{pd}_R(M)+\operatorname{depth}(M).
\]
\end{proposition}
\begin{proof}
We prove this by induction on $\operatorname{depth}(M)$. The most interesting case is the case $\operatorname{depth}(M)=0$. In this case, let
\[
0\to R^{n_e}\to R^{n_{e-1}}\to\ldots\to R^{n_0}\to M\to0
\]
be a minimal finite free resolution, so $e=\operatorname{pd}_R(M)$. By Lemma 10.102.2 we may assume all matrix coefficients of the maps in the complex are contained in the maximal ideal of $R$. Then on the one hand, by Proposition 10.102.9 we see that $\operatorname{depth}(R)\geq e$. On the other hand, breaking the long exact sequence into short exact sequences
\begin{align*}
0\to R^{n_e}\to R^{n_{e-1}}\to K_{e-2}\to0,\\
0\to K_{e-2}\to R^{n_{e-2}}\to K_{e-3}\to0,\\
\ldots,\\
0\to K_0\to R^{n_0}\to M\to0
\end{align*}
we see, using Lemma 10.72.6, that
\begin{align*}
\operatorname{depth}(K_{e-2})&\geq\operatorname{depth}(R)-1,\\
\operatorname{depth}(K_{e-3})&\geq\operatorname{depth}(R)-2,\\
&\ldots,\\
\operatorname{depth}(K_0)&\geq\operatorname{depth}(R)-(e-1),\\
\operatorname{depth}(M)&\geq\operatorname{depth}(R)-e
\end{align*}
and since $\operatorname{depth}(M)=0$ we conclude $\operatorname{depth}(R)\leq e$. This finishes the proof of the case $\operatorname{depth}(M)=0$.

Induction step. If $\operatorname{depth}(M)>0$, then we pick $x\in\mathfrak m$ which is a nonzerodivisor on both $M$ and $R$. This is possible, because either $\operatorname{pd}_R(M)>0$ and $\operatorname{depth}(R)>0$ by the aforementioned Proposition 10.102.9 or $\operatorname{pd}_R(M)=0$ in which case $M$ is finite free hence also $\operatorname{depth}(R)=\operatorname{depth}(M)>0$. Thus $\operatorname{depth}(R\oplus M)>0$ by Lemma 10.72.6 (for example) and we can find an $x\in\mathfrak m$ which is a nonzerodivisor on both $R$ and $M$. Let
\[
0\to R^{n_e}\to R^{n_{e-1}}\to\ldots\to R^{n_0}\to M\to0
\]
be a minimal resolution as above. An application of the snake lemma shows that
\[
0\to(R/xR)^{n_e}\to(R/xR)^{n_{e-1}}\to\ldots\to(R/xR)^{n_0}\to M/xM\to0
\]
is a minimal resolution too. Thus $\operatorname{pd}_R(M)=\operatorname{pd}_{R/xR}(M/xM)$. By Lemma 10.72.7 we have $\operatorname{depth}(R/xR)=\operatorname{depth}(R)-1$ and $\operatorname{depth}(M/xM)=\operatorname{depth}(M)-1$. Till now depths have all been depths as $R$ modules, but we observe that $\operatorname{depth}_R(M/xM)=\operatorname{depth}_{R/xR}(M/xM)$ and similarly for $R/xR$. By induction hypothesis we see that the Auslander-Buchsbaum formula holds for $M/xM$ over $R/xR$. Since the depths of both $R/xR$ and $M/xM$ have decreased by one and the projective dimension has not changed we conclude.
\end{proof}


\subsection*{整理说明（非作者正文）}
本条只计 Auslander--Buchsbaum 公式一个问题。正文先列作者的支撑材料，再列主命题及完整证明；各段内部保留原文顺序，没有将所附引理和正合性判据另作题数。将部分长复形改为折行公式，\verb|\text| 中的算子改作 \verb|\operatorname|；不改变数学内容。

标准先修保留原文引用：深度引理（10.72.6）、模去一个非零因子的深度变化（10.72.7）、蛇引理、局部环上有限投射模自由、伴随素理想和素理想回避。正合性判据使用的 10.102.3、10.102.6、10.102.7 分别是深度零时有限自由正合复形的分裂、分裂复形的秩计算、对非零因子取商后截断复形的正合性；不再递归重证这些先修。

原主证明的 $\operatorname{depth}(R)\geq e$ 并未只留作一个未展开的引用：所用 10.102.9 的正合性与正则序列联系已附完整作者证明。在线正文中的简略指标写法和推导表达照录；网页评论不作为作者证明。没有取得整个原生章节文件，保存的是实际使用的网页 TeX 正文摘录。

\subsection*{可修改的简短标注（非作者正文）}
$c$：有限生成模的深度与投射维数。

$a$：极小自由消解；子式理想；正则序列；伴随素理想处的局部化；深度引理；非零因子取商；蛇引理。

\texttt{cross\_theory\_status = yes}。以自由消解的长度表示同调维数，通过子式理想中的正则序列控制局部环的深度，再以非零因子取商同步降低两个深度量而保持消解长度，得到等式。原文位置：10.102.9 的正合性判据及 10.111.1 的深度零情形和归纳步骤。全部标注为非穷尽、可修改初稿。
\subsection*{许可与修改记录}
原数学正文作者：The Stacks Project Authors（集体作者）。原文按 GNU Free Documentation License 1.2 或后续版本发布；本摘录的整理部分随原文采用相同许可。许可全文随包保存在 \texttt{sources/stacks/GFDL-1.2.txt}。无不变章节、封面文字或封底文字。本文题名与来源作品题名不同，不暗示原作者认可本次选题或标注。

\textbf{History.} 2026-09-28：ChatGPT 为本对话材料包整理摘录，加入题号、来源定位、独立排版与简短标注；数学正文的具体排版变动见整理说明。

\end{document}
